\SourceLeaf{340}
and hence
\[
\begin{array}{rl}
(4,1) : (4,\text{III}) & = (3,1) : (\text{I}, \text{III}) \\
(4,1) : (4,\text{II}) & = (2,1) : (\text{I}, \text{II}) \\
(4,3) : (4,\text{II}) & = (2,3) : (\text{III}, \text{II}).
\end{array}
\]
Dividing the last two equations gives
\begin{equation}
(4,1) : (4,3) = \frac{(2,1)}{(2,3)} : \frac{(\text{I}, \text{II})}{(\text{III}, \text{II})}
= \frac{(2,1)}{(3,2)} : \frac{(\text{I}, \text{II})}{(\text{II}, \text{III})}.
\tag{1}
\end{equation}
If we add to the system of three equations
\[
(4,2) : (4,\text{III}) = (3,2) : (\text{II}, \text{III}),
\]
we obtain
\begin{equation}
(4,1) : (4,2) = \frac{(3,1)}{(3,2)} : \frac{(\text{I}, \text{III})}{(\text{II}, \text{III})}
\tag{2}
\end{equation}
and similarly
\begin{equation}
(4,2) : (4,3) = \frac{(2,1)}{(3,1)} : \frac{(\text{I}, \text{II})}{(\text{I}, \text{III})}.
\tag{3}
\end{equation}
From (1) we obtain
\[
\frac{(\text{I}, \text{II})}{(\text{II}, \text{III})} = \frac{(2,1) . (4,3)}{(4,1) . (3,2)},
\]
and from (2)
\[
\frac{(\text{I}, \text{III})}{(\text{II}, \text{III})} = \frac{(3,1) . (4,2)}{(4,1) . (3,2)}.
\]
Consequently
\[
(\text{I}, \text{II}) : (\text{II}, \text{III}) : (\text{III}, \text{I}) = (2,1) . (4,3) :
(4,1) . (3,2) : (4,2) . (3,1).
\]
If we denote the angles as before and set
\[
\sphericalangle\ \text{I}, \text{II}, \text{III} = [\text{I}, \text{III}],
\]
we have
\[
(\text{I}, \text{II}) : (\text{II}, \text{III}) : (\text{III}, \text{I}) = \sin[\text{I}, \text{II}] :
\sin[\text{II}, \text{III}] : \sin[\text{III}, \text{I}].
\]
Now
\[
\begin{array}{rl}
{[\text{I}, \text{II}]} & = [\text{I}, 4] + [4, \text{II}] = \alpha + \delta \\
{[\text{II}, \text{III}]} & = [\text{II}, 4] + [4, \text{III}] = \gamma + \zeta \\
{[\text{III}, \text{I}]} & = [\text{III}, 4] + [4, \text{I}] = \beta + \varepsilon,
\end{array}
\]
hence
\[
\frac{\sin(\alpha+\delta)}{(2,1) . (4,3)} = \frac{\sin(\gamma+\zeta)}{(4,1) . (3,2)} =
\frac{\sin(\beta+\varepsilon)}{(4,2) . (3,1)},
\]
which agrees with the earlier equation.

The problem takes a similar form if 4 lies outside triangle 1, 2, 3. We then take any
one of the four vertices, for\,example 4 again, draw (4,1), (4,2), (4,3) and determine the points
I, II, III in the same way and then the triangle (I), (II), (III), again seeking the relations between
sides and angles.
